\documentclass[11pt]{article}
\usepackage{fontspec,amsmath,amssymb,amsthm,hyperref}
\usepackage[margin=1in]{geometry}
\setlength{\emergencystretch}{3em}
\title{Zero-dimensional parameter and state spaces}
\author{LangevinLean mathematical exposition}
\date{}
\begin{document}\maketitle
\noindent Source module: \texttt{IPLA/Degenerate.lean}. Compile this proof document with XeLaTeX or LuaLaTeX. The accompanying HTML pairs the mathematical argument with the exact Lean statements and records the scope of the formalization.


Euclidean space over an empty coordinate type contains a single vector. These elementary cases allow the public theorem to retain independent dimensions that may both be zero.


\section{The parameter error vanishes when d\ensuremath{\theta} = 0}


If \(d_\theta=0\), every parameter vector equals every other one. In particular \(p_\theta z=\theta_*\), so \(\|p_\theta z-\theta_*\|^2=0\) pointwise. Its integral under any measure, and the square root of that integral, are zero. No probability, potential or noise assumption is required.




\textbf{Correspondence with the displayed declarations.}


\begin{itemize}

\item \texttt{parameterRMS\_zero\_dimension}: There is no parameter error when the parameter dimension is zero. This holds for every latent dimension and every measure.
\end{itemize}


\noindent Lean declarations: \texttt{parameterRMS\_zero\_dimension}.

\section{Every term of the full zero-dimensional bound vanishes}


If also \(d_x=0\), the whole state is a singleton. Any two probability measures on it are equal and their \(W_2\) distance is zero. The equilibrium radius \(\sqrt{2d_\theta/(Nm)}\) and discretization factor \(\sqrt{d_x+d_\theta/N}\) are also zero. Substituting these facts reduces the full claimed inequality to \(0\le0\), irrespective of the values of \(m,\gamma,L\). This handles the case in which the nonempty-coordinate analytic theorem cannot be instantiated.




When \(d_\theta=0\), the parameter space is a singleton, hence every parameter error vector and its square vanish. If also \(d_x=0\), the whole state space is a singleton. Any probability law there is the same Dirac law, so its Wasserstein distance to the Gibbs probability is zero. The dimension factors in the remaining terms are also zero, giving the public bound in this case directly, without invoking a theorem requiring a nonempty coordinate index.




\textbf{Correspondence with the displayed declarations.}


\begin{itemize}

\item \texttt{zero\_dimensions\_bound}: The complete displayed error bound in the zero-dimensional state. No regularity of a potential, noise law, or initialization is needed here.
\end{itemize}


\noindent Lean declarations: \texttt{zero\_dimensions\_bound}.

\end{document}
